%=================================================================
\documentclass[education,article,submit,moreauthors]{Definitions/mdpi}

% Journal short name for Education Sciences: education
\firstpage{1}
\makeatletter
\setcounter{page}{\@firstpage}
\makeatother
\pubvolume{1}
\issuenum{1}
\articlenumber{0}
\pubyear{2026}
\copyrightyear{2026}
\datereceived{ }
\daterevised{ }
\dateaccepted{ }
\datepublished{ }
%\pdfoutput=1 % Uncomment for arXiv upload

\usepackage{seqsplit}
\newcommand{\Pten}{\ensuremath{\mathrm{Precision}@10\%}}
\newcommand{\Rten}{\ensuremath{\mathrm{Recall}@10\%}}
\newcommand{\AP}{\ensuremath{\mathrm{AP}}}
\newcommand{\AUROC}{\ensuremath{\mathrm{AUROC}}}

\Title{Evaluating Data-Minimized Early-Warning Systems for Educator Decision Support: Evidence Across Review Capacities and Educational Stages}

\newcommand{\orcidauthorA}{0000-0001-6006-4813}
\newcommand{\orcidauthorB}{0000-0003-2331-6326}
\newcommand{\orcidauthorC}{0000-0003-4615-5756}
\newcommand{\orcidauthorD}{0009-0004-2567-173X}
\newcommand{\orcidauthorE}{0000-0002-8330-4282}
\newcommand{\orcidauthorF}{0009-0006-5319-7742}

\Author{Aigul Shaikhanova $^{1,2}$\orcidA{}, Oleksandr Kuznetsov $^{3,4,*}$\orcidB{}, Gulmira Shangytbayeva $^{5,*}$\orcidC{}, Kamila Bakenova $^{1,6}$\orcidD{}, Kainizhamal Iklassova $^{7}$\orcidE{} and Dana Tulemisova $^{8}$\orcidF{}}

\AuthorNames{Aigul Shaikhanova, Oleksandr Kuznetsov, Gulmira Shangytbayeva, Kamila Bakenova, Kainizhamal Iklassova and Dana Tulemisova}

\address{%
$^{1}$ \quad Global Education and Training, iSchool at Illinois, University of Illinois at Urbana-Champaign, Champaign, IL 61820, USA\\
$^{2}$ \quad Department of Information Security, L.N. Gumilyov Eurasian National University, Satpayev 2, 010008 Astana, Kazakhstan\\
$^{3}$ \quad Department of Theoretical and Applied Sciences, eCampus University, Via Isimbardi 10, 22060 Novedrate (CO), Italy\\
$^{4}$ \quad Department of Intelligent Software Systems and Technologies, School of Computer Science and Artificial Intelligence, V.~N.~Karazin Kharkiv National University, 4 Svobody Sq., 61022 Kharkiv, Ukraine\\
$^{5}$ \quad Department of Informatics and Information Technologies, Zhubanov University, 34 Aliya Moldagulova Avenue, 030000 Aktobe, Kazakhstan\\
$^{6}$ \quad Center for Scientific Research and Commercialization, Arkalyk Pedagogical University, 17 Auelbekov str., 110300 Arkalyk, Kazakhstan\\
$^{7}$ \quad Department of Information and Communication Technologies, Manash Kozybayev North Kazakhstan University, Pushkin Str.~86, 150000 Petropavlovsk, Kazakhstan\\
$^{8}$ \quad School of Cybersecurity, Astana IT University, 55/11 Mangilik El Avenue, 010000 Astana, Kazakhstan}

\corres{Correspondence: oleksandr.kuznetsov@uniecampus.it (O.K.); gshangytbayeva@zhubanov.edu.kz (G.S.)}

\abstract{Educational early-warning systems can help educators decide which students may need closer attention when time and support capacity are limited. Institutions also need to know whether collecting less background information preserves the practical usefulness of these systems. We evaluated this question in the Open University Learning Analytics Dataset (OULAD) and two separately analyzed UCI education datasets. Models using all available fields were compared with models that omitted selected student-background information at several educational stages and at review capacities of 5\%, 10\%, 20\%, and 30\%. The primary outcome was the yield of each review list: the proportion of flagged students who later experienced the adverse educational outcome. Paired bootstrap intervals accounted for repeated comparisons across stages and capacities. Removing sex/gender information did not produce a family-wise difference in review-list yield in any dataset. Removing socioeconomic and family-related information had more substantial consequences: global predictive performance decreased throughout OULAD and UCI 697, with additional losses in selected larger OULAD review lists and at UCI 697 enrollment. Results for the smaller UCI 320 samples were inconclusive. Separate diagnostic models also showed that excluded attributes remained partly predictable from other records; direct deletion therefore did not remove all associated information. These post-development findings do not estimate the effect of alerts on teaching or student outcomes. They instead provide an evaluation framework for responsible integration: institutions should assess early-warning systems at the stage, staffing capacity, and data scope in which educators will actually use them, and should prospectively evaluate workload, professional judgement, and student support after deployment.}

\keyword{learning analytics; educational early-warning systems; educator decision support; data minimization; review capacity; technology integration; responsible educational technology}

\begin{document}

\section{Introduction}

Educational early-warning systems are intended to help educators notice students who may benefit from timely support. In practice, however, a prediction is useful only when it can be incorporated into professional work. A teacher, adviser, or student-support team may be able to review only a small proportion of a cohort, and the system consequently helps determine whose situation is examined first. Its educational value therefore depends not only on overall predictive accuracy, but also on the quality of the review list produced at the institution's actual staffing capacity.

This distinction has a direct practical meaning. If an institution can review 10\% of a cohort, review-list precision is the proportion of students in that 10\% list who later experience the adverse outcome; recall is the proportion of all subsequent adverse outcomes represented in the list. Precision--recall measures are particularly informative when outcomes are imbalanced and capacity is constrained \citep{saito2015precision}. Yet a useful retrospective list is only a prerequisite for effective support. Whether an alert improves counselling, teaching, retention, or achievement must ultimately be established in prospective educational practice \citep{plak2022early}.

The choice of information used to construct a review list is equally important. Student-background variables may improve prediction, but their collection raises questions about necessity, proportionality, interpretation, and governance. Data minimization is a lifecycle design principle; deleting a field does not by itself make a system fair, private, or lawful \citep{gdpr2016,drachsler2016delicate,rets2023six}. Information related to an excluded attribute may remain in other records, and two models with similar aggregate performance may still direct educators' attention to different students. Evaluating data minimization therefore requires the intended educational stage, staff capacity, comparison model, and consequences for review-list membership to be stated explicitly.

Previous work has established the importance of capacity-based alerting, calibration, subgroup auditing, repeated prediction stages, and alert-list stability \citep{wu2026calibrated,wuluo2026zenodo}. Leakage-aware OULAD protocols likewise show why only information genuinely available at the prediction time should be used \citep{le2026leap}. Privacy-conscious engagement-only prediction and causal or adversarial analyses of demographic information further demonstrate that predictive utility, privacy, and fairness are related but distinct questions \citep{chiu2026privacy,hasan2022utility}. What remains less clear is whether a conclusion reached at one review capacity or educational stage remains valid when the working conditions of educators change.

Within the broader study of effective digital-technology integration into educators' professional activity, this research develops and evaluates a capacity-aware framework for assessing early-warning systems before deployment. We compare review lists across four realistic capacities, examine several points in the educational process, distinguish targeted removal of sex/gender information from broader removal of socioeconomic and family-related information, and test whether excluded attributes remain statistically recoverable from retained records. The OULAD data and several elements of the temporal modeling design were used in an earlier article by overlapping investigators \citep{shaikhanova2025interpretable}; the present contribution is the systematic evaluation of data-removal conclusions across capacities, stages, and datasets. OULAD, UCI 697, and UCI 320 address different educational settings and are therefore interpreted separately rather than pooled as replications.

The study addresses four questions:

\begin{enumerate}[leftmargin=1.8em]
\item How does removing student-background information change the usefulness of review lists when educators can examine 5\%, 10\%, 20\%, or 30\% of a cohort?
\item Do these consequences change as additional evidence becomes available during the course or program?
\item Do targeted sex/gender removal and broader socioeconomic/family removal have different practical consequences across the three educational datasets?
\item Can an excluded attribute still be predicted from retained records, and what does this imply for responsible institutional use?
\end{enumerate}

\section{Materials and Methods}

\subsection{Study design, evidentiary status, and provenance}

This is a retrospective methodological evaluation of early-warning models using existing, de-identified educational datasets. The framework was developed iteratively from earlier analyses and then locked before the final multi-capacity, attribute-predictability, and harmonized-block evaluations reported here. No review capacity, educational stage, feature definition, model grid, alert rule, resampling plan, or reporting threshold was changed after that lock. The study should therefore be read as a rigorous post-development stress test rather than as preregistered confirmation.

A provenance audit identified that an earlier recovery copy of OULAD encoded 1,111 missing area-deprivation values differently from the official archive. Because this changed preprocessing, all current OULAD results were regenerated from the official pinned bytes and compared with controls from the same execution. The original release was also replayed exactly on the official archive. Full checksums, the transition between runs, and the computational evidence ledger are provided in the Supplementary Materials and public reproducibility package.

Figure~\ref{fig:studyframework} summarizes the study as an educator-facing decision process. The datasets supply evidence available at different educational stages, the paired models test alternative information policies, and the capacity rule translates model scores into the number of students a responsible team can realistically review.

\begin{figure}[H]
\centering
\includegraphics[width=\textwidth]{fig0_capacity_aware_framework.pdf}
\caption{Capacity-aware framework for evaluating an educational early-warning system before prospective use. OULAD, UCI 697, and UCI 320 represent different educational settings and were analyzed separately. The framework links the intended educational decision and evidence available at that stage to paired model comparisons, realistic staff capacity, and measures of practical usefulness.}
\label{fig:studyframework}
\end{figure}

\subsection{OULAD setting, outcome, and temporal partitions}

OULAD contains 32,593 student--module--presentation registrations across 22 presentations of seven anonymized modules, together with assessment and virtual-learning-environment records \citep{kuzilek2017oulad}. At each observation day, the analysis included registrations that had begun and had not already ended in withdrawal. The educational outcome distinguished future Fail or Withdrawn from Pass or Distinction. Data from 2013B and 2013J were used for fitting, 2014B for model selection and calibration, and 2014J for temporal evaluation. Review lists were constructed after days 14, 28, 42, 56, 70, and 90.

Because students who withdraw leave later risk sets, the size and composition of the evaluated cohort change over time. Comparisons across observation days describe how available evidence and review lists change; they do not identify the optimal time for an educational intervention. A sensitivity analysis followed registrations present across the first four or all six observation days, while recognizing that this conditions on later continuation.

\subsection{Time-aware evidence and data-removal scenarios}

Only information available by an observation day was used. Virtual-learning-environment features summarized cumulative and recent activity, active days, timing, inactivity, pre-start activity, and change in click activity. An assessment submission counted only when both its due date and submission date were no later than the observation day. Scores were excluded because OULAD does not record when grades became visible to educators or students.

All standard models used course context, previous attempts, studied credits, registration timing, and the activity and assessment-process evidence available at that stage. The \emph{full-information} model additionally used gender, region, highest education, Index of Multiple Deprivation band, age band, and disability. The \emph{broadly minimized} model omitted all six fields, whereas the \emph{targeted} model omitted only gender and disability. Additional comparisons removed one background field at a time. The term \emph{data-minimized} describes these model inputs; it does not imply that the source dataset or wider institutional lifecycle no longer contains the fields.

\subsection{Predictive models, selection, and calibration}

Histogram gradient boosting was implemented in scikit-learn \citep{pedregosa2011scikit}. Missing numeric values were median-imputed with missingness indicators; categorical variables were most-frequent-imputed and ordinal-encoded. Every transformation was fitted on training data only. Model complexity was selected on the validation presentation by average precision, with simpler models preferred when validation results were tied. A positive-slope Platt model was then fitted on validation predictions and transferred unchanged to the temporal evaluation presentation. Calibration was evaluated rather than assumed \citep{vancalster2019calibration,collins2024tripod,moons2025probast}. Complete hyperparameters are reported in the Supplementary Materials.

All comparisons pair a data-minimized model with a full-information control created under the same data split, feature-availability rules, model family, and software environment. Selected models and predictions were fixed before the uncertainty calculations.

\subsection{Review capacity and educational interpretation of the metrics}

Within each operational group $g$---a module in OULAD, a course in UCI 697, or a school in UCI 320---students were ordered by estimated risk. At review capacity $b$, the number placed in the review list was
\[
  k_g=\left\lceil b n_g\right\rceil, \qquad b\in\{0.05,0.10,0.20,0.30\}.
\]
The primary measure was review-list precision, the share of reviewed students who later experienced the adverse outcome. Review-list recall measured the share of all later adverse outcomes included in that list:
\[
  P_b=\frac{\sum_g TP_g}{\sum_g k_g}, \qquad
  R_b=\frac{\sum_g TP_g}{\sum_g N_{+,g}}.
\]
For example, \Pten{} is the proportion of students in the highest-risk 10\% review list who later experienced the outcome, while \Rten{} is the proportion of all such outcomes captured in that list. Exact score ties at a list boundary were resolved by a prespecified deterministic key; a sensitivity analysis averaged over possible orders within tied boundary scores.

Average precision (AP) summarized ranking quality across possible thresholds, area under the receiver operating characteristic curve (AUROC) measured overall discrimination, and the Brier score assessed probability accuracy. The original OULAD decision rule treated a loss greater than three percentage points in day-42 \Pten{} as practically concerning. Because this reference was chosen during development, it is used as an operational guide rather than a formal non-inferiority margin.

\subsection{OULAD paired uncertainty and common-cohort sensitivity}

Uncertainty was estimated by resampling students, retaining all registrations belonging to a resampled student, and reconstructing the full-information and data-minimized review lists in every replicate. The paired design asks how the two policies differ for the same resampled educational records. The day-42 analysis used 5,000 replicates and the other OULAD stages used 2,000. Family-wise intervals accounted for the six observation days within each prespecified policy family. Models were not refitted during resampling.

The four-horizon common cohort contained 8,853 registrations from 8,536 students; the six-horizon cohort contained 8,507 registrations from 8,216 students. Within each, day-42-minus-day-14 difference-in-differences compared the change in each exclusion gap with the full-control change. These analyses condition on later registration retention.

\subsection{Fallback models without clickstream data}

To examine a setting in which learning-platform activity is unavailable, day-42 fallback models excluded all clickstream variables and retained course context and assessment-process information that was due and observable by that day. An enriched version added further process counts and ratios, but no grades, click telemetry, or future records. Model selection and calibration followed the same validation procedure as the standard models.

At a 10\% review capacity, the prespecified operational rule allowed the enriched fallback to be no more than five percentage points below the standard broadly minimized comparator. This comparison asks whether a clickstream-free fallback remains practically adequate, not merely whether it improves on a weaker fallback.

\subsection{External educational datasets and cross-fitted controls}

Two official UCI datasets were evaluated after OULAD and are reported as distinct external analyses. UCI 697, \emph{Predict Students' Dropout and Academic Success}, contains 4,424 higher-education students, 36 predictors, and the outcomes Dropout, Enrolled, and Graduate \citep{realinho2021uci697}. The primary comparison distinguished Dropout from Graduate and excluded students still enrolled; a secondary outcome compared Dropout with all other statuses. Enrollment models used only information available on entry. Semester-one models added first-semester curricular-unit evidence while prohibiting all second-semester fields. Data-removal scenarios omitted direct personal information, family and financial information, or gender and educational special-needs status. A sparse operational comparator retained only a small set of admissions and course variables.

UCI 320, \emph{Student Performance}, provides separate Mathematics ($n=395$) and Portuguese ($n=649$) samples from two secondary schools \citep{cortez2008uci320}. The subjects were not pooled because the supplied linkage fields did not uniquely identify the same students across files. The outcome was final-grade failure ($G3<10$). Baseline models excluded all grades; the next stages added G1 and then G2 while always excluding the final grade G3. Data-removal scenarios omitted sex, demographic and family fields, sensitive and wellbeing fields, or retained only school-observable information.

Both UCI analyses used ten shared stratified outer folds. Model selection and calibration were conducted only within each outer training fold, and logistic regression served as a sensitivity model. Review capacity was allocated within course for UCI 697 and within school for UCI 320. Paired resampling used 5,000 replicates for the primary gradient-boosting comparisons and 2,000 for secondary comparisons. Family-wise intervals covered the two UCI 697 or three UCI 320 educational stages within each policy family. These datasets broaden the settings examined, but they are not treated as independent confirmations because their analysis followed the OULAD work.

\subsection{Evaluation across four review capacities}

Before computing the final capacity analysis, review fractions of 5\%, 10\%, 20\%, and 30\% were fixed. For each fraction we reconstructed the number of reviewed students, true and false alerts, precision, recall, lift over prevalence, overlap between review lists, and the proportion of students retained from the full-information list. Existing OULAD test probabilities and UCI out-of-fold probabilities were reused, so this step changed the decision capacity but did not alter or refit the models.

\subsection{Residual predictability after direct attribute removal}

Separate diagnostic models asked whether an excluded attribute could still be predicted from the information retained in the corresponding data-minimized model. Logistic regression was the primary diagnostic, with gradient boosting as a nonlinear sensitivity. OULAD diagnostics targeted gender and disability and were evaluated temporally on 2014J, including a subset of students not observed during fitting. UCI 697 diagnostics targeted gender and educational special-needs status, and UCI 320 diagnostics targeted sex within the saved outer folds. Educational outcomes, identifiers, split fields, future-stage information, and the target attribute itself were prohibited. These diagnostics measure residual predictability: they do not show that an early-warning model used a particular proxy or that a fairness or privacy violation occurred.

\subsection{Comparable information groups across datasets}

Two information groups were defined to support transparent comparisons across datasets. The \emph{sex/gender} group comprised the corresponding recorded field in each dataset. The \emph{socioeconomic/family} group comprised highest education and area deprivation in OULAD; parental qualifications and occupations, debtor and tuition status, and scholarship status in UCI 697; and parental education and occupations in UCI 320. These are operationally comparable groups, not claims that the datasets measure an identical social construct.

Existing paired comparisons were reused where available. Only missing gradient-boosting comparisons needed for the two information groups were fitted, each with a full-information control created in the same execution. Controls were required to reproduce the corresponding locked probabilities to numerical precision before a comparison was accepted. No UCI 697 outcome model was refitted. A previously preserved XuetangX analysis was excluded because the official dataset schema did not support verifiable demographic comparisons.

\subsection{Cross-dataset paired uncertainty and multiplicity}

Review lists were reconstructed within every bootstrap replicate. OULAD resampled students and retained repeated registrations, whereas the UCI analyses resampled rows. The same resampling multiplicities were applied to full-information and data-minimized predictions at all four capacities. Primary gradient-boosting comparisons used 5,000 replicates and sensitivity comparisons used 2,000. For review-list metrics, each family covered all prespecified stages and capacities for one dataset, outcome, model, and data-removal policy; global probability metrics covered stages. We report pointwise and Bonferroni-adjusted family-wise 95\% percentile intervals. The primary comparison was the difference in review-list precision between the data-minimized and full-information models. Recall, list overlap, AP, AUROC, and Brier-score differences were secondary. An interval that includes zero is interpreted as inconclusive, not as proof of equivalence.

\subsection{Computational reproducibility}

Inputs, data partitions, feature-availability rules, fitted pipelines, predictions, resampling draws, tables, and figures were checked by executable verification scripts. The final audit independently reconstructed all primary capacity results and statistical intervals from preserved artifacts. Complete counts, hashes, and replay tolerances are reported in the Supplementary Materials and aggregate-only public repository. These checks establish internal reproducibility; they are not a substitute for evaluation in another institution.

\section{Results}

\subsection{The provenance-corrected OULAD analysis provided a stable basis for comparison}

The corrected analysis used the official OULAD archive and an explicit rule for tied scores. Small intermediate counts changed relative to the superseded recovery run, but no result was adjusted toward a historical value. The exact transition is documented in the Supplementary Materials so that the educational comparisons below rest on one consistent data source.

At day 42, 9,040 active registrations were eligible for review and 3,584 later ended in Fail or Withdrawn. At a 10\% capacity, each model produced a list of 908 registrations. The full-information model identified 769 subsequent adverse outcomes, the broadly minimized model identified 762, and the targeted model without gender and disability identified 765.

\subsection{Removing all six background fields had stage-dependent consequences}

At day 42, the broadly minimized model achieved a review-list precision of 0.8392, compared with 0.8469 for the full-information model. In practical terms, both lists contained about 84 subsequent adverse outcomes for every 100 registrations reviewed. The paired difference was $-0.0077$ (95\% interval $[-0.0265,0.0155]$), remaining within the three-percentage-point operational reference used in the original analysis.

The same conclusion did not hold at every stage. At day 14, the broadly minimized list identified 670 subsequent adverse outcomes among 935 reviewed registrations, compared with 708 for the full-information model. The precision difference was $-0.0406$, and the family-wise interval remained below zero ($[-0.0692,-0.0036]$). Average precision was also lower at every observation day. Broad removal of background information was therefore not uniformly cost-free: its practical effect depended on when the list was produced.

\subsection{Targeted removal largely preserved the yield of educator review lists}

At day 42, removing gender and disability produced a review-list precision of 0.8425, compared with 0.8469 for the full-information model: a difference of less than half a percentage point (Table~\ref{tab:day42}). The interval included zero. Small AUROC and Brier-score differences nevertheless favored the full-information model, showing that a similar yield at one review capacity does not mean that the models rank all students or estimate probabilities identically.

\begin{table}[H]
\caption{Day-42 comparison of targeted gender/disability removal with the full-information model.}
\label{tab:day42}
\begin{tabularx}{\textwidth}{LRRRR}
\toprule
\textbf{Metric} & \textbf{G/D excluded} & \textbf{Full} & \textbf{Difference} & \textbf{95\% interval} \\
\midrule
Alerts & 908 & 908 & --- & --- \\
True-risk alerts & 765 & 769 & $-4$ & --- \\
False alerts & 143 & 139 & $+4$ & --- \\
Precision@10\% & 0.8425 & 0.8469 & -0.0044 & [-0.0110, 0.0099] \\
Recall@10\% & 0.2134 & 0.2146 & -0.0011 & [-0.0028, 0.0025] \\
Average precision & 0.6934 & 0.6943 & -0.0009 & [-0.0022, 0.0005] \\
AUROC & 0.7585 & 0.7603 & -0.0018 & [-0.0027, -0.0009] \\
Brier score & 0.1915 & 0.1909 & +0.0006 & [0.0003, 0.0008] \\
\bottomrule

\end{tabularx}
\noindent{\footnotesize{Intervals use 5,000 paired student-level bootstrap replicates and are conditional on the fitted models. Positive Brier differences favor the full-information model.}}
\end{table}

Across the six observation days, precision differences ranged from $-0.0044$ at day 42 to $+0.0035$ at day 90 (Table~\ref{tab:landmarks}; Figure~\ref{fig:partial}). Every pointwise interval included zero, and all family-wise lower bounds remained above the three-percentage-point reference. The targeted policy therefore showed no practically large loss in review-list yield under the evaluated conditions, although this evidence is not a formal equivalence or non-inferiority test.

\begin{table}[H]
\centering
\caption{Targeted gender/disability removal versus the full-information model at a 10\% review capacity.}
\label{tab:landmarks}
\scriptsize
\setlength{\tabcolsep}{3.5pt}
\begin{tabular}{rrrrrrrrr}
\toprule
Day & Reviewed & Targeted TP & Targeted P & Full TP & Full P & Difference & Pointwise 95\% & Family-wise 95\% \\
\midrule
14 & 935 & 709 & 0.7583 & 708 & 0.7572 & +0.0011 & [-0.0107, 0.0128] & [-0.0149, 0.0171] \\
28 & 923 & 763 & 0.8267 & 764 & 0.8277 & -0.0011 & [-0.0108, 0.0141] & [-0.0173, 0.0195] \\
42 & 908 & 765 & 0.8425 & 769 & 0.8469 & -0.0044 & [-0.0110, 0.0099] & [-0.0143, 0.0134] \\
56 & 890 & 803 & 0.9022 & 801 & 0.9000 & +0.0022 & [-0.0067, 0.0124] & [-0.0101, 0.0165] \\
70 & 873 & 813 & 0.9313 & 815 & 0.9336 & -0.0023 & [-0.0092, 0.0069] & [-0.0115, 0.0092] \\
90 & 855 & 804 & 0.9404 & 801 & 0.9368 & +0.0035 & [-0.0082, 0.0070] & [-0.0117, 0.0093] \\
\bottomrule

\end{tabular}
\end{table}

\begin{figure}[H]
\centering
\includegraphics[width=0.96\textwidth]{fig1_partial_six_horizons.png}
\caption{Difference in review-list precision after targeted removal of gender and disability. Dark intervals are pointwise paired bootstrap intervals; pale intervals account for the six observation days. The dashed line marks the three-percentage-point operational reference.}
\label{fig:partial}
\end{figure}

This finding was specific to the capacity-constrained review list. At day 14, overall ranking quality measured by average precision decreased by $-0.00342$ (family-wise interval $[-0.00613,-0.00071]$). At day 42, the small AUROC difference ($-0.0018$) and Brier-score difference ($+0.0006$) also favored the full-information model. Thus, similar performance for the first 10\% of students did not imply identical behavior across the full cohort.

The students brought to educators' attention also changed. Compared with the full-information list, targeted removal replaced between 17 and 61 registrations in each direction, depending on the observation day. Broad removal replaced 290 registrations at day 14 and 169 at day 42. Institutions should therefore examine who enters and leaves a review list even when its overall yield changes little.

\subsection{Earlier stages were more sensitive to prior-education information}

The largest single-field loss occurred when highest education was removed at day 14: review-list precision decreased by $-0.0342$ (95\% interval $[-0.0544,-0.0085]$). By day 42, all six single-field differences lay between $-0.0066$ and $+0.0077$. This pattern suggests that prior educational information may matter most before richer evidence of current participation becomes available, but it does not establish that the field itself causes the outcome.

Following the same registrations across observation days did not establish a consistent improvement in the targeted policy relative to the full-information model. The broader removal policy showed more favorable point estimates at later stages, but family-wise intervals included zero. Later evidence can reduce dependence on enrollment information, yet these retrospective comparisons cannot determine whether waiting produces a better intervention opportunity.

\subsection{Clickstream-free fallback models improved but remained substantially less useful}

Adding richer assessment-process information improved the clickstream-free fallback: average precision increased by 0.0256, AUROC by 0.0315, and the Brier score improved by 0.0070 relative to the sparse fallback. Its 10\% review-list precision increased by only $+0.0055$, with an interval that included zero. Point estimates are shown in Table~\ref{tab:strict}.

\begin{table}[H]
\centering
\caption{Day-42 results for models that do not use clickstream information.}
\label{tab:strict}
\footnotesize
\setlength{\tabcolsep}{4pt}
\begin{tabular}{lrrrr}
\toprule
Policy & Model size & Relevant cases & \Pten & AP \\
\midrule
Six-field-excluded standard & 7 & 762 & 0.8392 & 0.6798 \\
Strict original & 31 & 669 & 0.7368 & 0.5883 \\
Strict enriched due-process & 7 & 674 & 0.7423 & 0.6139 \\
Strict typed alternative & 7 & 670 & 0.7379 & 0.5945 \\
\bottomrule

\end{tabular}
\noindent{\footnotesize{Model size is the selected number of terminal nodes.}}
\end{table}

Against the standard broadly minimized model, however, the enriched fallback achieved a review-list precision of 0.7423 rather than 0.8392: a loss of 9.69 percentage points (95\% interval $[-0.1289,-0.0738]$). It therefore did not meet the five-percentage-point operational rule. Removing clickstream dependence may be desirable for resilience or governance, but the resulting reduction in practical utility must be measured and an explicit fallback or abstention policy is needed.

\subsection{Early prediction depended strongly on family and financial information in UCI 697}

For the primary Dropout-versus-Graduate outcome, the full-information model achieved a 10\% review-list precision of 0.9111 at enrollment and 0.9865 after semester one. A sparse admissions-based model achieved 0.5849 and 0.9730, respectively (Table~\ref{tab:uci697}). Thus, at enrollment, reviewing 100 students selected by the sparse model would identify about 33 fewer subsequent dropout cases than reviewing 100 students selected by the full-information model. After first-semester evidence became available, the difference narrowed to 1.35 percentage points, with the family-wise interval reaching zero.

\begin{table}[H]
\centering
\caption{UCI 697 review-list precision for the primary terminal-status outcome.}
\label{tab:uci697}
\small
\begin{tabularx}{\textwidth}{LRRRR}
\toprule
Stage & Full information & Sparse operational & Difference (points) & Family-wise 95\% CI \\
\midrule
Enrollment & 0.9111 & 0.5849 & -32.61 & [-39.25, -27.06] \\
After semester 1 & 0.9865 & 0.9730 & -1.35 & [-4.04, 0.00] \\
\bottomrule

\end{tabularx}
\end{table}

The type of information removed mattered. At enrollment, removing family and financial fields reduced 10\% review-list precision by 21.83 percentage points. Removing seven direct personal fields changed it by only $-0.54$ points, and removing gender plus special-needs status changed it by $+0.81$ points; both intervals included zero. Small differences in overall ranking remained detectable, however, and after semester one some global metrics still favored the full-information model. Stronger current academic evidence narrowed the practical gap, but it did not make every reduced model interchangeable.

\subsection{Later grades improved prediction, but exclusion effects remained uncertain in UCI 320}

The UCI 320 full-information models improved markedly after grades became available (Figure~\ref{fig:external}). In Mathematics, average precision increased from 0.432 before grades to 0.825 with G1 and 0.941 with G1 and G2; 10\% review-list precision increased from 0.575 to 0.925 and 0.975. In Portuguese, the corresponding values increased from 0.412 to 0.663 and 0.845 for average precision, and from 0.409 to 0.591 and 0.727 for review-list precision. Complete point estimates are reported in the Supplementary Materials.

\begin{figure}[H]
\centering
\includegraphics[width=0.96\textwidth]{fig2_external_stage_contrasts.pdf}
\caption{How later academic evidence changed prediction in two distinct datasets. Panel A shows UCI 697 review-list precision for full-information and sparse operational models. Panel B shows UCI 320 average precision as G1 and G2 become available; Mathematics and Portuguese are analyzed separately. The panels are not pooled estimates or direct replications.}
\label{fig:external}
\end{figure}

By contrast, the consequences of removing background fields were uncertain. Nearly all pointwise and family-wise intervals included zero for ranking, probability accuracy, and review-list precision. The only exception was a small Mathematics period-one AUROC improvement after removing sex. Given the small samples and the many settings examined, this isolated result does not support a general policy conclusion. UCI 320 demonstrates the informational value of staged grades, while also showing why a data-removal effect observed in one educational setting should not be assumed elsewhere.

\subsection{Review capacity changed the practical effect of data minimization}

The practical conclusion changed with the number of students an institution could review. In OULAD, broad data minimization produced lower point estimates in 19 of the 24 stage--capacity combinations. After family-wise correction, losses remained at the 30\% capacity on days 14, 28, and 56 (Table~\ref{tab:c14budget}). Targeted removal of gender and disability produced no family-wise difference at any evaluated stage or capacity. A conclusion based only on the highest-risk 10\% therefore did not fully describe what happened when educators worked further down the ranked list.

In UCI 697, the sparse operational model remained substantially below the full-information model at enrollment at all four capacities, with losses ranging from 34.92 percentage points at 5\% to 17.97 points at 30\%. After semester one, differences remained detectable only at the larger 20\% and 30\% capacities. Removing family and financial fields showed the same pattern. No UCI 320 review-list comparison remained different after family-wise correction. Later academic evidence reduced some information gaps, but the degree of reduction depended on how many students educators were expected to review.

\begin{table}[H]
\centering
\caption{Selected changes in review-list precision after data minimization at four institutional review capacities. Values are percentage points; asterisks identify family-wise intervals that exclude zero.}
\label{tab:c14budget}
\scriptsize
\setlength{\tabcolsep}{4pt}
\begin{tabular}{lrrrr}
\toprule
Educational setting & 5\% & 10\% & 20\% & 30\% \\
\midrule
OULAD day14, minimized & -3.62 & -4.06 & -2.19 & -2.07$^{*}$ \\
OULAD day28, minimized & +0.22 & -0.22 & -1.52 & -1.99$^{*}$ \\
OULAD day42, minimized & +1.75 & -0.77 & -0.50 & -1.51 \\
OULAD day56, minimized & -0.22 & +0.34 & -0.96 & -1.99$^{*}$ \\
OULAD day70, minimized & +0.68 & -0.46 & -0.69 & -1.53 \\
OULAD day90, minimized & -0.23 & +0.94 & -0.82 & -0.43 \\
UCI 697 enrollment, legacy & -34.92$^{*}$ & -32.61$^{*}$ & -22.37$^{*}$ & -17.97$^{*}$ \\
UCI 697 semester1, legacy & -3.17 & -1.35 & -3.55$^{*}$ & -3.56$^{*}$ \\
\bottomrule

\end{tabular}
\end{table}

\subsection{Deleted attributes remained partly predictable from other records}

Every primary diagnostic model discriminated the excluded attribute above chance in held-out data (Table~\ref{tab:c14proxy}; Figure~\ref{fig:c14proxy}). OULAD gender AUROC ranged from 0.690 to 0.812 across observation days, and disability ranged from 0.663 to 0.680. Sex AUROC was approximately 0.75 in UCI 320 Mathematics and 0.73 in Portuguese; UCI 697 gender reached 0.771--0.774. Educational special-needs status in UCI 697 was much less stable because only 40 positive cases were available and average precision was approximately 0.020 at a prevalence of 0.011.

The association remained when OULAD evaluation was limited to students not observed during diagnostic-model fitting. The retained records therefore contained information associated with attributes that had been removed from the early-warning model. This finding does not identify the source of the association or show that the outcome model relied on it. Its practical meaning is narrower but important: direct deletion should be treated as one governance measure, not as evidence that all related information has disappeared.

\begin{table}[H]
\centering
\caption{Held-out predictability of directly excluded attributes. Ranges cover the evaluated stages; AP lift is average precision divided by target prevalence.}
\label{tab:c14proxy}
\scriptsize
\begin{tabular}{llrrrr}
\toprule
Dataset & Target & Stages & Positive cases & AUROC range & AP-lift range \\
\midrule
OULAD & disability & 6 & 761--868 & 0.663--0.680 & 1.52--1.68 \\
OULAD & gender & 6 & 4905--5425 & 0.690--0.812 & 1.31--1.43 \\
UCI 320 Math. & sex & 3 & 187 & 0.746--0.750 & 1.50--1.50 \\
UCI 320 Port. & sex & 3 & 266 & 0.734--0.734 & 1.63--1.64 \\
UCI 697 & Educational special needs & 2 & 40 & 0.637--0.655 & 1.77--1.88 \\
UCI 697 & Gender & 2 & 1249 & 0.771--0.774 & 1.89--1.90 \\
\bottomrule

\end{tabular}
\end{table}

\begin{figure}[H]
\centering
\includegraphics[width=0.88\textwidth]{fig4_c14_proxy_persistence.pdf}
\caption{Held-out AUROC ranges after direct attribute removal. Each segment spans the evaluated stages and the point is their mean; the dashed line represents chance discrimination. Residual predictability does not show that the early-warning model used a specific proxy.}
\label{fig:c14proxy}
\end{figure}

\subsection{Sex/gender and socioeconomic/family information had different consequences}

Across 30 paired comparisons, removing sex/gender information produced no family-wise difference in review-list precision in OULAD, UCI 697, or UCI 320 (Table~\ref{tab:c14harmonized}). This should not be read as complete equality: small differences remained in selected global metrics, and the students included in the lists were not always identical.

Socioeconomic and family-related information showed a different pattern. In OULAD, removing it reduced overall ranking and probability accuracy at every observation day; review-list precision also decreased at the 30\% capacity on days 28 and 56. In UCI 697, removal produced a large enrollment-time loss and smaller remaining losses after semester one. UCI 320 again provided no stable review-list difference. Figure~\ref{fig:c14budget} shows the OULAD and UCI 697 capacity profiles separately.

\begin{table}[H]
\centering
\caption{Review-list precision after removal of comparable information groups across stages and capacities. Ranges are point estimates in percentage points; the final columns count family-wise intervals that remain below or above zero.}
\label{tab:c14harmonized}
\scriptsize
\begin{tabular}{llrrrr}
\toprule
Setting & Information removed & Comparisons & Point range & Below zero & Above zero \\
\midrule
OULAD & sex/\allowbreak{}gender & 24 & [-0.23, +1.54] & 0 & 0 \\
OULAD & socioeconomic/\allowbreak{}family & 24 & [-4.05, +0.94] & 2 & 0 \\
UCI 320 Math. & sex/\allowbreak{}gender & 12 & [-5.00, +5.00] & 0 & 0 \\
UCI 320 Math. & socioeconomic/\allowbreak{}family & 12 & [-10.00, +5.04] & 0 & 0 \\
UCI 320 Port. & sex/\allowbreak{}gender & 12 & [-4.55, +2.94] & 0 & 0 \\
UCI 320 Port. & socioeconomic/\allowbreak{}family & 12 & [-4.55, +8.82] & 0 & 0 \\
UCI 697/\allowbreak{}Dropout vs all other & socioeconomic/\allowbreak{}family & 8 & [-21.78, -0.87] & 7 & 0 \\
UCI 697/\allowbreak{}Dropout vs graduate & sex/\allowbreak{}gender & 8 & [-1.09, +0.81] & 0 & 0 \\
UCI 697/\allowbreak{}Dropout vs graduate & socioeconomic/\allowbreak{}family & 8 & [-23.81, -1.35] & 6 & 0 \\
\bottomrule

\end{tabular}
\end{table}

\begin{figure}[H]
\centering
\includegraphics[width=\textwidth]{fig3_c14_budget_harmonized.pdf}
\caption{Change in review-list precision after removing socioeconomic and family-related information. Open black circles identify family-wise intervals that exclude zero. Lines connect capacities within one educational stage. OULAD and UCI 697 use different populations, outcomes, and field definitions and are not pooled.}
\label{fig:c14budget}
\end{figure}

\subsection{The strongest losses were concentrated rather than universal}

Across 296 planned stage--capacity results, 30 family-wise intervals were below zero and none was above zero. Some rows describe the same paired comparison under two reporting views and are not independent confirmations. The central finding is therefore not that data minimization usually reduces review-list precision. Rather, the strongest losses were concentrated in larger OULAD review lists and in UCI 697 models that omitted family and financial information or used only sparse admissions data. Targeted sex/gender removal and all UCI 320 review-list comparisons remained inconclusive after family-wise correction.

Figure~\ref{fig:resultspractice} brings the principal findings together in the form most relevant to institutional decision making: what was observed, how cautiously it can be interpreted, and what an educator-facing implementation should do next.

\begin{figure}[H]
\centering
\includegraphics[width=\textwidth]{fig5_results_to_practice.pdf}
\caption{From empirical findings to responsible educational practice. The diagram summarizes the interpretation supported by the retrospective analyses and the corresponding institutional response. The arrows organize the decision pathway; they do not represent causal effects.}
\label{fig:resultspractice}
\end{figure}

\section{Discussion}

\subsection{Principal interpretation}

The central contribution is a practical evaluation framework for data-minimized early-warning systems. A statement such as ``removing background data has little cost'' is incomplete unless it specifies which information was removed, when the prediction was made, how many students educators could review, and which aspect of performance was considered. In OULAD, broad removal appeared acceptable at the original day-42 10\% operating point but produced detectable losses in several 30\% review lists. Targeted removal of gender and disability did not produce a family-wise loss in review-list yield, whereas removing socioeconomic and family-related information reduced overall predictive performance throughout OULAD and affected selected larger lists.

The external datasets clarify why this is an educational rather than purely technical issue. In UCI 697, sparse admissions data and the omission of family and financial information substantially weakened enrollment-time prioritization. First-semester achievement evidence reduced these gaps, but waiting for stronger evidence may also reduce the time available for support. UCI 320 showed pronounced gains after grades became available but did not yield a stable data-removal effect. The appropriate stage is therefore not simply the one with the highest predictive accuracy; it is the stage that balances reliable prioritization with a meaningful opportunity to act.

Review-list yield and whole-cohort performance answer different questions. A data-minimized model may identify a similar number of subsequent adverse outcomes within a fixed staff capacity while still changing the order of the remaining cohort, the estimated probabilities, or the students selected. Conversely, a model can lose some global discrimination without changing the yield at one particular capacity. Educational institutions should therefore evaluate both the usefulness of the actual review list and the distribution of attention it creates.

Residual attribute predictability adds a distinct governance lesson. Gender, disability, and special-needs information remained partly recoverable from other educational records even after direct removal. This is consistent with overlapping information and correlated social or behavioral processes, but it does not establish that the outcome model used a specific proxy. Direct deletion alone neither removes statistical association nor demonstrates fairness \citep{gardner2019fairness,pham2025fairness}. A stronger fairness claim would require subgroup allocation and performance analyses, a model-specific investigation, and an explicit normative standard.

The clickstream-free results illustrate the same need for context. Assessment-process evidence improved a sparse fallback but did not recover the practical yield of the standard comparator. Institutions may still need a fallback when platform telemetry is unavailable or should not be used, but the system should then disclose reduced confidence, abstain where appropriate, and avoid presenting a degraded list as operationally equivalent.

\subsection{Why provenance matters for educational decisions}

The provenance correction is not merely a software detail. A different representation of missing area-deprivation values changed how the full-information and partially minimized models interpreted student records. An implicit rule for tied scores also changed two selections at a review-list boundary. In a teacher-facing system, such changes can determine which individual students appear for review. Reproducibility must therefore cover source data, missing-value semantics, information timing, list construction, and software versions, not only the final model.

\subsection{Implications for educator-facing digital technology}

Module- or course-specific lists approximate the work faced by an educator or local support team more closely than one institution-wide ranking. Their precision, false-alert burden, membership changes, and fallback performance can show whether a digital tool is ready for prospective evaluation. Large-scale teacher-facing implementations and human-centred dashboard research nevertheless demonstrate that workflow, explanation, and professional interpretation must be studied with educators \citep{herodotou2019teachers,alfredo2024dashboard}. Data-driven decision making should treat platform evidence as one input to contextual professional judgement rather than as an automated instruction \citep{evenstein2026trusting}.

The results support a staged institutional procedure. First, define the educational decision, the responsible professional, and the point at which support remains possible. Second, state how many students can realistically be reviewed. Third, compare full-information and data-minimized systems at that capacity, examining list yield, false alerts, probability quality, list membership, and subgroup allocation. Fourth, repeat the evaluation when new academic evidence becomes available. Fifth, specify what the system does when important data are delayed, missing, or unavailable. A teacher-facing interface should explain why a student entered the review list, allow professional override, and record both the action taken and its outcome.

This procedure connects the study to the quality of educators' professional activity without replacing educational judgement with an algorithm. The present analysis evaluates the reliability and governance of a digital decision-support instrument. A subsequent prospective phase should examine whether the instrument reduces unproductive workload, improves the consistency and timeliness of professional decisions, supports appropriate interventions, and benefits students.

\subsection{Limitations}

The analysis was developed iteratively and is not preregistered confirmation. Its bootstrap intervals are conditional on fitted models, selected hyperparameters, and the positive-slope Platt calibration used here; they do not include model-selection, calibration-method, institutional-transport, or development-process uncertainty. A monotone calibration map preserves score ordering, but an isotonic analysis could introduce ties and alter probability-quality measures. The primary model family was gradient boosting, with logistic regression used only as a sensitivity in the UCI analyses. Family-wise intervals cover the prespecified stages and capacities within each named policy and metric, not every analysis conducted during the broader project.

The OULAD outcome combines Fail and Withdrawn, although these events may require different educational responses and occur on different timelines. Later risk sets exclude earlier withdrawals, and the common-cohort analyses condition on later continuation. The available timestamps also do not reconstruct the exact moment when every assessment result became visible. Consequently, the study describes prediction at defined data snapshots; it does not identify an optimal intervention time or distinguish intervention pathways for failure and withdrawal.

The three datasets differ in platform, educational level, outcome, timing, and sample size. UCI 320 contains small subject-specific samples, and its Mathematics and Portuguese files could not be uniquely linked or pooled. The common socioeconomic/family label also covers different observed fields in each dataset and should not be interpreted as measurement invariance. These differences are useful for testing the portability of conclusions, but they prevent a pooled effect or a claim of external replication.

Residual predictability is descriptive and does not establish actual proxy use, privacy leakage, or unfair allocation by the early-warning models. The special-needs diagnostic in UCI 697 is particularly uncertain because only 40 positive cases were available. We did not conduct a complete subgroup fairness audit, privacy attack, or prospective evaluation of educator workload, intervention uptake, possible harms, or student benefit. Those analyses require an institutionally defined use case and participation by the educators and students affected.

Finally, the official XuetangX schema contained no demographic fields, and a preserved derivative could not be independently connected to that schema. It was therefore excluded from all demographic evidence rather than used as unverifiable support. Detailed provenance exclusions, superseded results, and reproducibility records remain available in the Supplementary Materials.

\section{Conclusion}

Educational early-warning systems allocate scarce professional attention as well as predict risk. This study shows why data minimization must be evaluated within that educational function. The consequences of removing information changed with the type of field, the point in the educational process, and the proportion of students that staff could realistically review. Targeted removal of sex/gender information did not produce a family-wise loss in review-list yield in the three datasets, while removing socioeconomic and family-related information reduced overall predictive performance in OULAD and UCI 697 and weakened selected review lists. The smaller UCI 320 analyses remained inconclusive, and excluded attributes remained partly predictable from other records.

For educational institutions, the practical recommendation is clear: define the educator's decision and available review capacity before choosing or simplifying a model; evaluate the resulting list at each intended stage; examine false alerts, calibration, list membership, and subgroup allocation; and provide an explicit fallback when data are unavailable. The system should support transparent professional judgement rather than automate consequential action. This capacity-aware framework provides a reproducible basis for selecting and governing digital decision-support tools. The next step is prospective evaluation with educators to determine whether the tool improves workload, decision quality, timeliness of support, and student outcomes in practice.

\appendixstart
\appendix

\section{Practical Interpretation Guide}

Table~\ref{tab:interpretation} translates the principal evaluation quantities into the questions faced by educators and educational institutions. The quantities should be interpreted together; no single metric establishes that an early-warning system is ready for deployment.

\begin{table}[H]
\centering
\caption{Plain-language interpretation of the evaluation quantities.}
\label{tab:interpretation}
\small
\begin{tabularx}{\textwidth}{p{0.18\textwidth}p{0.34\textwidth}X}
\toprule
Quantity & Plain-language meaning & Institutional question \\
\midrule
Review capacity & The proportion of students that educators can realistically examine & How many cases can the responsible team review and follow up? \\
Precision at capacity $b$ & Among students placed in the review list, the proportion who later experience the adverse outcome & How much of the review effort is directed toward students represented by the outcome definition? \\
Recall at capacity $b$ & The proportion of all later adverse outcomes included in the review list & How many subsequent cases are missed because the list is limited? \\
AP and AUROC & How well the model ranks students across the full range of possible thresholds & Does the model retain useful ordering beyond the chosen review boundary? \\
Brier score & How closely predicted probabilities correspond to observed outcomes & Can the stated risk values be interpreted as probabilities? \\
List overlap & The extent to which two models bring the same students to educators' attention & Does data minimization change who is reviewed even when aggregate yield is similar? \\
\bottomrule
\end{tabularx}
\end{table}

\vspace{6pt}

\supplementary{The following supporting information can be downloaded at \linksupplementary{s1}: Table S1: Computational evidence and provenance ledger; Table S2: Review-list precision families across stages and capacities; Table S3: Comparable information-group contrasts at a 10\% review capacity; Table S4: Complete residual-attribute predictability results; Table S5: OULAD main-policy point estimates; Tables S6--S7: UCI review-capacity contrasts; Table S8: Single-field removal results; Table S9: Common-cohort timing sensitivity; Table S10: Review-list overlap; Tables S11--S12: Provenance-transition records; Table S13: Production and verification inventory. Supplementary Materials S1 is provided as a separate MDPI supplementary-file document and is numbered independently from the main text.}

\authorcontributions{Conceptualization, O.K.; methodology, O.K.; software, O.K., K.B. and K.I.; validation, O.K., K.B. and K.I.; formal analysis, A.S.; investigation, G.S.; resources, K.B. and K.I.; data curation, K.B. and K.I.; writing---original draft preparation, O.K.; writing---review and editing, D.T.; visualization, D.T.; supervision, A.S.; project administration, G.S.; funding acquisition, A.S. All authors have read and agreed to the published version of the manuscript.}

\funding{This research was funded by the Science Committee of the Ministry of Science and Higher Education of the Republic of Kazakhstan, grant number AP23489228.}

\institutionalreview{Not applicable. This study used only publicly accessible, de-identified secondary datasets (OULAD; UCI ``Predict Students' Dropout and Academic Success''; UCI ``Student Performance''), with no recruitment, contact, intervention, re-identification attempt, or access to non-public records.}

\informedconsent{Not applicable. This study used only pre-existing, publicly available, de-identified secondary data; no human subjects were recruited or contacted by the authors.}

\dataavailability{OULAD is available from the UCI Machine Learning Repository at \url{https://archive.ics.uci.edu/dataset/349/open+university+learning+analytics+dataset} and is described by \citet{kuzilek2017oulad}. The official archive used here had SHA-256 \texttt{\seqsplit{f2ed1902616c1fe8d2824d872c0b7d2d72be435bf0124d077044fe4be2c6d3e4}}. UCI 697 is available at \url{https://archive.ics.uci.edu/dataset/697/predict+students+dropout+and+academic+success}; the outer ZIP used here had SHA-256 \texttt{\seqsplit{e90e55fd65ec462ae283ebeb2cca409319e3460ed898d8754f62fb35cc83a65d}}. UCI 320 is available at \url{https://archive.ics.uci.edu/dataset/320/student+performance}; the outer ZIP had SHA-256 \texttt{\seqsplit{82ae9d66437b9808df42e8c89d2bb179c46e9cfbcf06f38abc1d20b3b747e177}}. The current nested \texttt{student.zip}, not \texttt{.student.zip\_old}, was analyzed. All three UCI pages identify the datasets as CC~BY~4.0. The aggregate-only reproducibility release accompanying this submission is available at \url{https://github.com/KuznetsovKarazin/SecureEWS}; it contains analysis code, the locked protocol, machine-readable aggregate results and bootstrap draws, verification reports, and manuscript sources. Raw or processed student-level data, individual predictions, fitted model bundles, private clean-room archives, and XuetangX derivative files are not redistributed.}

\acknowledgments{During the preparation of this manuscript, the authors used OpenAI ChatGPT and Codex to assist with code development, verification workflows, visualization, and language and drafting support. The authors reviewed and validated the code, analyses, references, figures, and final text, revised the outputs as necessary, and take full responsibility for the content of this publication.}

\conflictsofinterest{The authors declare no conflicts of interest. The funders had no role in the design of the study; in the collection, analyses, or interpretation of data; in the writing of the manuscript; or in the decision to publish the results.}

%\isPreprints{}{% This command is only used for ``preprints''.
\begin{adjustwidth}{-\extralength}{0cm}
%} % If the paper is ``preprints'', please uncomment this parenthesis.

\reftitle{References}

\bibliography{references}

\end{adjustwidth}

\end{document}
